\documentclass[10pt,a4paper]{amsart}
\usepackage{amsmath,amssymb,amsthm,graphicx}
% Retained minimal preamble for the Stacks Project chapter spaces-descent.tex (Descent on Spaces).
% Verbatim excerpt of the chapter preamble (sources/stacks-spaces-descent/preamble.tex): the reference/slogan/history
% environment block (native lines 7-12), Xy-pic setup (native lines 14-19) and the native theorem declarations with
% the mathematical macros (native lines 154-217). Copied mechanically; nothing added, removed or redefined. The
% wrapper class and the amsmath/amssymb/amsthm/graphicx base packages are supplied by the assembler, so the
% class/xr-hyper/hyperref/externaldocument part of the native preamble is omitted.
% For dealing with references we use the comment environment
\usepackage{verbatim}
\newenvironment{reference}{\comment}{\endcomment}
%\newenvironment{reference}{}{}
\newenvironment{slogan}{\comment}{\endcomment}
\newenvironment{history}{\comment}{\endcomment}
% For commutative diagrams we use Xy-pic
\usepackage[all]{xy}

% We use 2cell for 2-commutative diagrams.
\xyoption{2cell}
\UseAllTwocells
% Theorem environments.
%
\theoremstyle{plain}
\newtheorem{theorem}[subsection]{Theorem}
\newtheorem{proposition}[subsection]{Proposition}
\newtheorem{lemma}[subsection]{Lemma}

\theoremstyle{definition}
\newtheorem{definition}[subsection]{Definition}
\newtheorem{example}[subsection]{Example}
\newtheorem{exercise}[subsection]{Exercise}
\newtheorem{situation}[subsection]{Situation}

\theoremstyle{remark}
\newtheorem{remark}[subsection]{Remark}
\newtheorem{remarks}[subsection]{Remarks}

\numberwithin{equation}{subsection}

% Macros
%
\def\lim{\mathop{\mathrm{lim}}\nolimits}
\def\colim{\mathop{\mathrm{colim}}\nolimits}
\def\Spec{\mathop{\mathrm{Spec}}}
\def\Hom{\mathop{\mathrm{Hom}}\nolimits}
\def\Ext{\mathop{\mathrm{Ext}}\nolimits}
\def\SheafHom{\mathop{\mathcal{H}\!\mathit{om}}\nolimits}
\def\SheafExt{\mathop{\mathcal{E}\!\mathit{xt}}\nolimits}
\def\Sch{\mathit{Sch}}
\def\Mor{\mathop{\mathrm{Mor}}\nolimits}
\def\Ob{\mathop{\mathrm{Ob}}\nolimits}
\def\Sh{\mathop{\mathit{Sh}}\nolimits}
\def\NL{\mathop{N\!L}\nolimits}
\def\CH{\mathop{\mathrm{CH}}\nolimits}
\def\proetale{{pro\text{-}\acute{e}tale}}
\def\etale{{\acute{e}tale}}
\def\QCoh{\mathit{QCoh}}
\def\Ker{\mathop{\mathrm{Ker}}}
\def\Im{\mathop{\mathrm{Im}}}
\def\Coker{\mathop{\mathrm{Coker}}}
\def\Coim{\mathop{\mathrm{Coim}}}

% Boxtimes
%
\DeclareMathSymbol{\boxtimes}{\mathbin}{AMSa}{"02}

%
% Macros for moduli stacks/spaces
%
\def\QCohstack{\mathcal{QC}\!\mathit{oh}}
\def\Cohstack{\mathcal{C}\!\mathit{oh}}
\def\Spacesstack{\mathcal{S}\!\mathit{paces}}
\def\Quotfunctor{\mathrm{Quot}}
\def\Hilbfunctor{\mathrm{Hilb}}
\def\Curvesstack{\mathcal{C}\!\mathit{urves}}
\def\Polarizedstack{\mathcal{P}\!\mathit{olarized}}
\def\Complexesstack{\mathcal{C}\!\mathit{omplexes}}
% \Pic is the operator that assigns to X its picard group, usage \Pic(X)
% \Picardstack_{X/B} denotes the Picard stack of X over B
% \Picardfunctor_{X/B} denotes the Picard functor of X over B
\def\Pic{\mathop{\mathrm{Pic}}\nolimits}
\def\Picardstack{\mathcal{P}\!\mathit{ic}}
\def\Picardfunctor{\mathrm{Pic}}
\def\Deformationcategory{\mathcal{D}\!\mathit{ef}}
\usepackage{cleveref}
\crefname{originallemma}{original lemma}{original lemmas}
\Crefname{originallemma}{Original lemma}{Original lemmas}
\makeatletter
\newlabel{bootstrap-lemma-descend-algebraic-space}{{0ADV}{source}{}{sourcepointer-1}{}}
\newlabel{bootstrap-lemma-descend-algebraic-space@cref}{{[originallemma][0][]0ADV}{[1][0][]source}{}{}{}}
\makeatother
\crefname{originallemma}{original lemma}{original lemmas}
\Crefname{originallemma}{Original lemma}{Original lemmas}
\makeatletter
\newlabel{spaces-topologies-definition-fppf-covering}{{03Y8}{source}{}{sourcepointer-2}{}}
\newlabel{spaces-topologies-definition-fppf-covering@cref}{{[originallemma][0][]03Y8}{[1][0][]source}{}{}{}}
\makeatother
\title{Descent data versus sheaves under an fppf covering of algebraic spaces: for an fppf covering \{X\_i -> X\} of algebraic spaces over S, descent data relative to the covering are equivalent to sheaves F on (Sch/S)\_fppf with a map F -> X such that every X\_i x\_X F is an algebraic space; F is an algebraic space iff the corresponding datum is effective.}
\author{The Stacks Project Authors}
\date{stacks-project:spaces-descent.tex:local-snapshot-20261012; GNU FDL 1.2 or later}
\begin{document}
\maketitle
\noindent\textit{Editorial note.} Statement span is lines 3803-3835 of sources/stacks-spaces-descent/pinned-original-source.tex; proof span is lines 3837-3881, the complete author proof (43 body lines).\par\medskip
\noindent\textit{Editorial note.} Navigation stubs: 2 cited-but-not-included labels are supplied as non-load-bearing pointers with their pinned tags where the snapshot tag index has them.\par\medskip
\noindent\textit{Editorial note.} extra\_packages is empty on purpose: the retained preamble already loads xy with options and a second declaration clashes in the assembler wrapper.\par\medskip
\noindent\textit{Editorial note.} Text is retained verbatim, including the author's own typos and placeholders.\par\medskip
\noindent\textit{Editorial note.} Pool stacks-spaces-descent, tightened-ruling accept list (construction + required verification in span; only the mutual-inverse check left to the reader).\par\medskip
\noindent\textit{Editorial note.} Span computed from the \textbackslash{}label marker to \textbackslash{}end\{lemma\}; precard\_gate combined span RESULT=OK (probe L9999).\par\medskip
\noindent\textit{Editorial note.} rqual sv-mismatch fix 2026-10-14: source\_version aligned to SOURCE.json (semicolon+git form, intake authority; package evidence.json already correct); revision bumped for re-assembly per publication rotation (takeover-line defer ticket, relayed by the s15 line).\par\medskip
\noindent\textit{Allowed standard prerequisites (editorial).}\begin{itemize}
\item Explicit prerequisite boundary. See why-field; every pointer named here occurs verbatim in the retained statement/proof text.
\end{itemize}\medskip
\ifdefined\hypertarget\hypertarget{sourcepointer-1}{}\fi
\noindent\textit{Source reference.} Original lemma 0ADV: sources/stacks-spaces-descent/pinned-original-source.tex, source lines 3881--3881. Cited by the retained statement or proof but not part of the collected spans; supplied only as a navigation pointer. The argument does not depend on its content.\par\medskip
\ifdefined\hypertarget\hypertarget{sourcepointer-2}{}\fi
\noindent\textit{Source reference.} Original lemma 03Y8: sources/stacks-spaces-descent/pinned-original-source.tex, source lines 3881--3881. Cited by the retained statement or proof but not part of the collected spans; supplied only as a navigation pointer. The argument does not depend on its content.\par\medskip
\begin{lemma}
\label{lemma-descent-data-sheaves}
Let $S$ be a scheme. Let $\{X_i \to X\}_{i \in I}$ be an fppf
covering of algebraic spaces over $S$ (Topologies on Spaces,
Definition \ref{spaces-topologies-definition-fppf-covering}).
There is an equivalence of categories
$$
\left\{
\begin{matrix}
\text{descent data }(V_i, \varphi_{ij})\\
\text{relative to }\{X_i \to X\}
\end{matrix}
\right\}
\leftrightarrow
\left\{
\begin{matrix}
\text{sheaves }F\text{ on }(\Sch/S)_{fppf}\text{ endowed}\\
\text{with a map }F \to X\text{ such that each}\\
X_i \times_X F\text{ is an algebraic space}
\end{matrix}
\right\}.
$$
Moreover,
\begin{enumerate}
\item the algebraic space $X_i \times_X F$ on the right hand side
corresponds to $V_i$ on the left hand side, and
\item the sheaf $F$ is an algebraic space\footnote{We will see
later that this is always the case if $I$ is not too large, see
Bootstrap, Lemma \ref{bootstrap-lemma-descend-algebraic-space}.}
if and only if the
corresponding descent datum $(X_i, \varphi_{ij})$ is effective.
\end{enumerate}
\end{lemma}

\begin{proof}
Let us construct the functor from right to left.
Let $F \to X$ be a map of sheaves on $(\Sch/S)_{fppf}$ such that each
$V_i = X_i \times_X F$ is an algebraic space. We have the
projection $V_i \to X_i$.
Then both $V_i \times_X X_j$ and $X_i \times_X V_j$
represent the sheaf $X_i \times_X F \times_X X_j$
and hence we obtain an isomorphism
$$
\varphi_{ii'} : V_i \times_X X_j \to X_i \times_X V_j
$$
It is straightforward to see that the maps $\varphi_{ij}$
are morphisms over $X_i \times_X X_j$ and satisfy the
cocycle condition. The functor from right to left is given
by this construction $F \mapsto (V_i, \varphi_{ij})$.

\medskip\noindent
Let us construct a functor from left to right.
The isomorphisms $\varphi_{ij}$ give isomorphisms
$$
\varphi_{ij} : V_i \times_X X_j \longrightarrow X_i \times_X V_j
$$
over $X_i \times X_j$. Set $F$ equal to the coequalizer in the
following diagram
$$
\xymatrix{
\coprod_{i, i'} V_i \times_X X_j
\ar@<1ex>[rr]^-{\text{pr}_0}
\ar@<-1ex>[rr]_-{\text{pr}_1 \circ \varphi_{ij}}
& &
\coprod_i V_i \ar[r]
&
F
}
$$
The cocycle condition guarantees that $F$ comes with a map
$F \to X$ and that $X_i \times_X F$ is isomorphic to $V_i$.
The functor from left to right is given
by this construction $(V_i, \varphi_{ij}) \mapsto F$.

\medskip\noindent
We omit the verification that these constructions
are mutually quasi-inverse functors. The final statements
(1) and (2) follow from the constructions.
\end{proof}

\end{document}
